\begin{initiativkarte}{FSR I2E}

  % Bild (optional)
  \IfFileExists{bilder/fsr_i2e.jpg}{%
    \begin{center}
      \includegraphics[width=\linewidth,height=5cm,keepaspectratio]{bilder/fsr_i2e.jpg}
    \end{center}
    \vspace{0.4cm}
  }{}

  % Kurzbeschreibung
  \textit{\textcolor{hawgray}{Der Fachschaftsrat I2E vertritt die Interessen der Studierenden des Departments Informatik und Informations- und Elektrotechnik und unterstützt sie während des gesamten Studiums.}}

  \vspace{0.5cm}
  \hawrule

  % Ausführliche Beschreibung
  \textbf{\textcolor{hawblue}{Was machen wir?}}\\[0.3cm]

  Der Fachschaftsrat (FSR) wird jedes Semester von den Studierenden des Departments gewählt und bildet das Bindeglied zwischen Studierenden, Lehrenden und Hochschulverwaltung.

  Zu den wichtigsten Aufgaben gehören:
  \begin{itemize}
      \item Vertretung der Interessen der Studierenden gegenüber der Hochschulverwaltung und den Lehrenden.
      \item Organisation und Unterstützung von Veranstaltungen wie Fachvorträgen, Workshops, Informationsveranstaltungen und studentischen Events.
      \item Beratung und Unterstützung bei studienrelevanten Fragen und Problemen.
      \item Förderung des Austauschs und der Vernetzung innerhalb der Studierendenschaft.
      \item Mitarbeit in Gremien und Ausschüssen der Hochschule.
      \item Ausleihe von Technik für studentische Projekte und Veranstaltungen.
  \end{itemize}

  Der FSR ist eine gute Anlaufstelle für alle Fragen rund um das Studium sowie für Ideen und Anregungen zur Verbesserung des Hochschullebens.

  \vspace{0.5cm}

  % Schlagwörter
  \tag{Mitbestimmung}
  \tag{Beratung}
  \tag{Events}
  \tag{Studierendenvertretung}
  \tag{Netzwerk}

  \vspace{0.6cm}
  \hawrule

  % Informationen
  \infozeile{\faUsers}{Zielgruppe: Alle Studierenden des Departments Informatik und Informations- und Elektrotechnik}
  \infozeile{\faGraduationCap}{Semester: Ab dem 1. Semester}
  \infozeile{\faCalendarAlt}{Ort: Berliner Tor 7, Raum 0.06 / 0.07}
  \infozeile{\faEnvelope}{Kontakt: \href{mailto:fsr_ei@haw-hamburg.de}{fsr\_ei@haw-hamburg.de}}
  \infozeile{\faGlobe}{Website: \href{https://i2e.fsr.haw-hamburg.de/}{i2e.fsr.haw-hamburg.de}}

\end{initiativkarte}

